\documentclass[10pt,letterpaper]{article}
\usepackage{cornell-notes}

\title{Clause 31.05.02.08: Constructors And Destructor}
\author{}
\date{}
\setDocTitle{Clause 31.05.02.08: Constructors And Destructor}
\setDocSubtitle{C++ 2024}
\setDocOwner{C++ 2024 Cornell Notes}
\setDocDate{\today}
\setCornellCollection{C++ 2024 Cornell Notes}
\setCornellUnitType{Clause}
\setCornellUnitNumber{31.05.02.08}
\setCornellUnitTitle{Constructors And Destructor}
\hypersetup{pdftitle={Clause 31.05.02.08: Constructors And Destructor},pdfsubject={Cornell notes},pdfauthor={}}

\begin{document}
\maketitle
\begin{CornellOverviewBox}{Overview}
\textbf{Classification:} Standard library - Normative\\[2pt]
\textbf{Learning objective:} Understand the rules, terminology, and semantic consequences associated with Constructors and destructor.
\end{CornellOverviewBox}\begin{CornellSummaryBox}{Summary}
{\sffamily\bfseries\color{Accent} Summary}\par\vspace{3pt}Section 31.5.2.8 focuses on Constructors and destructor. Apply the specified interface together with its mandates, effects, result conditions, exception behavior, and complexity constraints. When applying it, preserve the distinction between mandatory requirements, recommendations, permissions, and behavior deliberately left undefined, unspecified, or implementation-defined.
\end{CornellSummaryBox}

\end{document}
